\documentclass[12pt]{article}

\usepackage{sbc-template}
\usepackage{graphicx,url}
\usepackage[utf8]{inputenc}
\usepackage[T1]{fontenc}
\usepackage[brazilian]{babel}
\usepackage{amsmath,amssymb}
\usepackage{booktabs}
\usepackage{textcomp}

% IPA symbols — use tipa but avoid the babel-active " character
\usepackage{tipa}
\newcommand{\ipaE}{\textipa{E}}%  open-mid front unrounded = ɛ
\newcommand{\ipaO}{\textipa{O}}%  open-mid back rounded = ɔ
\newcommand{\ipaG}{\textipa{G}}%  voiced velar fricative = ɣ
\newcommand{\ipag}{\textipa{g}}%  IPA g (U+0261)
\newcommand{\stressmark}{\raisebox{0.25ex}{\textquotesingle}}% ˈ (IPA stress)

\sloppy

\title{Distance-Aware Loss: Penalidade Fonologicamente Graduada\\para Convers\~ao Grafema-Fonema no Portugu\^es Brasileiro}

\author{[AN\^ONIMO]\inst{1}}

\address{[Afilia\c{c}\~ao anonimizada]
  \email{[email anonimizado]}
}

\begin{document}

\maketitle

\begin{abstract}
Este trabalho apresenta uma fun\c{c}\~ao de custo fonologicamente informada
--- a \emph{Distance-Aware} (DA) Loss --- para sistemas de convers\~ao
grafema-fonema (G2P) no Portugu\^es Brasileiro. Diferente da CrossEntropy
padr\~ao, que trata todos os erros de substitui\c{c}\~ao igualmente, a DA
Loss penaliza erros proporcionalmente \`a dist\^ancia articulat\'oria entre
o fonema predito e o alvo, ponderada pela confian\c{c}a do modelo.
A dist\^ancia \'e calculada sobre 24 features articulat\'orias do
PanPhon. Aplicada a um BiLSTM encoder-decoder com aten\c{c}\~ao de Bahdanau,
a DA Loss redistribui sistematicamente os erros das classes catastr\'oficas
(Classe~D) para classes mais pr\'oximas (Classe~B). O sistema \'e avaliado
sobre um conjunto de teste estratificado de 28.782 palavras
(${\sim}$181k fonemas) --- escala 57$\times$ maior que avalia\c{c}\~oes
compar\'aveis em PT-BR --- alcan\c{c}ando PER de 0,48\% (IC Wilson 95\%
[0,46\%; 0,51\%]) e WER de 5,33\% na configura\c{c}\~ao de refer\^encia;
uma configura\c{c}\~ao complementar sem separadores sil\'abicos atinge
WER de 4,96\%. Testes em 31 palavras fora do
vocabul\'ario mostram 100\% de acur\'acia em palavras reais in\'editas
do PT-BR, evid\^encia consistente com generaliza\c{c}\~ao de regras
fonol\'ogicas.
\end{abstract}

\begin{resumo}
  This work presents a phonologically-informed loss function ---
  Distance-Aware (DA) Loss --- for Grapheme-to-Phoneme (G2P) systems
  in Brazilian Portuguese. Unlike standard cross-entropy, DA~Loss
  penalizes substitution errors proportionally to the articulatory
  distance between predicted and target phonemes, weighted by the
  model's prediction confidence. Applied to a BiLSTM encoder-decoder,
  DA~Loss systematically redistributes errors from catastrophic
  (Class~D) to phonologically closer substitutions (Class~B). The
  system is evaluated on a stratified 28{,}782-word test set
  (${\sim}$181k phonemes), achieving 0.48\% PER (Wilson 95\% CI
  [0.46\%; 0.51\%]) and 5.33\% WER; a complementary configuration
  without syllable separators reaches 4.96\% WER. Generalization tests
  on 31 out-of-vocabulary words yield 100\% accuracy on genuine novel
  PT-BR words.
\end{resumo}

\section{Introdu\c{c}\~ao}

A convers\~ao autom\'atica de texto escrito em representa\c{c}\~ao
fon\'etica --- tarefa \emph{Grapheme-to-Phoneme} (G2P) --- \'e componente
fundamental de sistemas de s\'{\i}ntese de fala (TTS), reconhecimento de
fala e processamento de linguagem natural. Para o Portugu\^es Brasileiro,
a tarefa apresenta desafios espec\'{\i}ficos:
(i)~\textbf{ambiguidade graf\^emica} --- o grafema ``c'' realiza-se como
/k/ em ``cama'' mas /s/ em ``cena''; o ``r'' em coda realiza-se /x/ em
final de palavra mas /\ipaG{}/ antes de consoante vozeada;
(ii)~\textbf{neutraliza\c{c}\~ao voc\'alica} --- em s\'{\i}labas \'atonas,
/\ipaE{}/$\leftrightarrow$/e/ e /\ipaO{}/$\leftrightarrow$/o/
neutralizam-se; (iii)~\textbf{suprassegmentais} --- o acento t\^onico
(\textsuperscript{\textprimstress})\footnote{S\'{\i}mbolo IPA para acento prim\'ario,
representado como token separado na sa\'{\i}da do modelo.} e fronteiras
sil\'abicas (.) s\~ao tokens n\~ao-articulat\'orios que exigem tratamento
especial na fun\c{c}\~ao de custo.

Modelos seq2seq treinados com CrossEntropy (CE) tratam todos os erros
igualmente: predizer /\ipaE{}/ quando o correto \'e /e/ --- par
m\'{\i}nimo com 1 feature de diferen\c{c}a --- incorre na mesma
penalidade que predizer /k/ para /a/ --- erro de 8+ features. Essa
cegueira fonol\'ogica distorce o sinal de treinamento: o modelo aprende
\emph{que} errou, mas n\~ao \emph{o quanto} errou.

Este trabalho prop\~oe a \textbf{Distance-Aware (DA) Loss}, que adiciona
ao sinal de treinamento um termo proporcional \`a dist\^ancia
articulat\'oria entre o fonema predito e o alvo, ponderado pela
confian\c{c}a do modelo. Aplicamos essa fun\c{c}\~ao a um BiLSTM
encoder-decoder com aten\c{c}\~ao de Bahdanau
\cite{bahdanau2014neural} --- arquitetura deliberadamente escolhida para
isolar a contribui\c{c}\~ao da fun\c{c}\~ao de custo de novidades
arquiteturais.

\textbf{Compara\c{c}\~ao com LatPhon.} A refer\^encia mais pr\'oxima \'e
o LatPhon \cite{chary2025latphon}, um Transformer multil\'{\i}ngue de 4
camadas (7,5M par\^ametros, RoPE) que reporta PER de 0,86\% (IC Wilson
95\% [0,56\%; 1,16\%]) em ${\sim}$500 palavras PT-BR do mesmo dicion\'ario
fon\'etico. Nosso sistema alcan\c{c}a PER de 0,48\% (IC [0,46\%; 0,51\%])
em 28.782 palavras. Os intervalos de confian\c{c}a n\~ao se
sobrep\~oem --- o limite superior do nosso sistema (0,51\%) fica abaixo
do limite inferior do LatPhon (0,56\%), diferen\c{c}a estatisticamente
significativa a 95\% de confian\c{c}a.

\textbf{Contribui\c{c}\~oes}: (1) DA Loss: objetivo de treinamento
fonologicamente graduado combinando dist\^ancia PanPhon com confian\c{c}a
de predi\c{c}\~ao; (2) avalia\c{c}\~ao em larga escala com split
estratificado; (3) an\'alise fatorial de capacidade $\times$ separadores
$\times$ DA Loss; (4) taxonomia de qualidade de erros (Classes~A--D);
(5) avalia\c{c}\~ao de generaliza\c{c}\~ao OOV em 6 categorias.

\section{Dados e Protocolo de Avalia\c{c}\~ao}

\subsection{Corpus}

O corpus de treinamento consiste em \textbf{95.937 pares}
(palavra, transcri\c{c}\~ao IPA) de um dicion\'ario fon\'etico do
Portugu\^es Brasileiro. O charset de entrada cobre a--z (exceto k, w, y,
ausentes do dicion\'ario) mais diacr\'{\i}ticos portugueses. Palavras
contendo k, w ou y s\~ao tratadas como OOV de caractere. Previamente ao
treinamento, 10.252 entradas foram corrigidas para a distin\c{c}\~ao
ASCII-g (U+0067) vs.\ IPA-\ipag{} (U+0261), necess\'aria para
\emph{lookup} correto de features PanPhon.

\subsection{Split Estratificado}

O corpus \'e dividido \textbf{60/10/30} (treino/valida\c{c}\~ao/teste)
com estratifica\c{c}\~ao por tr\^es vari\'aveis fonol\'ogicas: tipo de
acento (ox\'{\i}tona / parox\'{\i}tona / proparox\'{\i}tona), faixa de
contagem de s\'{\i}labas (1, 2, 3, 4, 5+) e faixa de comprimento em
grafemas ($\leq$4, 5--7, 8--10, 11+). A combina\c{c}\~ao gera
${\sim}$48 estratos. Qualidade do split: $\chi^2{=}0{,}95$ ($p{=}0{,}678$),
Cram\'er $V{=}0{,}0007$ --- aus\^encia de diferen\c{c}a distribucional
significativa.

\textbf{Evid\^encia emp\'{\i}rica de vi\'es de split.} Um split
n\~ao-estratificado 70/10/20 (Exp0, mesma arquitetura 4,3M, CE) atingiu
PER de 1,12\%, enquanto o split estratificado 60/10/30 (Exp1) atingiu
0,66\% --- redu\c{c}\~ao de 41\% no PER atribu\'{\i}vel inteiramente ao
protocolo de avalia\c{c}\~ao \cite{kohavi1995crossvalidation}.

\subsection{M\'etricas}

\textbf{PER} (Phoneme Error Rate): dist\^ancia de Levenshtein sobre
sequ\^encias de fonemas, normalizada pelo comprimento da refer\^encia
\cite{morris2004cer,bisani2008joint}. \textbf{WER}: fra\c{c}\~ao de
palavras com qualquer erro. \textbf{IC Wilson 95\%}
\cite{wilson1927probable,brown2001interval}: para ${\sim}$181k fonemas,
IC de $\pm$0,03 p.p.

\textbf{M\'etricas graduadas}: \textbf{PER$_w$} (PER ponderado), onde
cada substitui\c{c}\~ao \'e ponderada pela dist\^ancia PanPhon. Classes
de erro: A (exato), B ($\leq$0,050, par m\'{\i}nimo), C ($\leq$0,150,
mesma fam\'{\i}lia), D ($>$0,150, classes diferentes).

\section{Arquitetura}

Utilizamos um \textbf{BiLSTM encoder-decoder com aten\c{c}\~ao de
Bahdanau} \cite{bahdanau2014neural}, arquitetura estabelecida para G2P
supervisionado \cite{rao2015g2p}.

\textbf{Encoder}: BiLSTM de 2 camadas. Para cada posi\c{c}\~ao $t$,
$h_t = [\overrightarrow{h_t}; \overleftarrow{h_t}]$ fornece contexto
bidirecional completo.

\textbf{Aten\c{c}\~ao}: a cada passo de decodifica\c{c}\~ao $t$:
\begin{equation}
e_{t,j} = v^\top \tanh(W_h h_j + W_s s_{t-1}), \quad
\alpha_{t,j} = \frac{\exp(e_{t,j})}{\sum_k \exp(e_{t,k})}, \quad
c_t = \sum_j \alpha_{t,j} h_j
\end{equation}

\textbf{Decoder}: LSTM de 2 camadas, \emph{teacher forcing} no treino,
autoregressivo na infer\^encia.

\textbf{Configura\c{c}\~oes testadas} (Tabela~\ref{tab:configs}).

\begin{table}[ht]
\centering
\caption{Configura\c{c}\~oes de arquitetura avaliadas.}
\label{tab:configs}
\begin{tabular}{lccc}
\toprule
\textbf{Config} & \textbf{Embedding} & \textbf{Hidden} & \textbf{Par\^ametros} \\
\midrule
Pequena       & 128D & 256D & 4,3M \\
Intermedi\'aria & 192D & 384D & 9,7M \\
Grande        & 256D & 512D & 17,2M \\
\bottomrule
\end{tabular}
\end{table}

\textbf{Embeddings fon\^emicos.}
Os experimentos principais utilizam embeddings aprendidos (inicializa\c{c}\~ao
Glorot). Adicionalmente, testamos embeddings inicializados com PanPhon
(24 features articulat\'orias projetadas para 128D), provendo um prior
geom\'etrico onde fonemas similares come\c{c}am pr\'oximos no espa\c{c}o
de embedding. Com CE, o PanPhon init produz erros qualitativamente
melhores (menor PER$_w$) mas n\~ao melhora o PER absoluto, pois a CE
n\~ao refor\c{c}a a estrutura articulat\'oria. A DA~Loss opera
independentemente da inicializa\c{c}\~ao do embedding: utiliza uma
tabela externa de dist\^ancias PanPhon, n\~ao a geometria do embedding.
Os dois mecanismos s\~ao ortogonais.

\subsection{Protocolo de Treinamento}

Todos os modelos s\~ao treinados com Adam ($\beta_1{=}0{,}9$,
$\beta_2{=}0{,}999$), taxa de aprendizado $10^{-3}$,
\emph{batch size}~96 e \emph{early stopping} com paci\^encia~10
sobre a loss de valida\c{c}\~ao. O \emph{teacher forcing ratio}
\'e 1,0 durante o treinamento.

O \emph{batch size} de 96 foi escolhido empiricamente: com 34 estratos
fonol\'ogicos no corpus, $\text{batch\_size} \approx 3 \times S$
garante $\geq$3 amostras por estrato por batch, reduzindo a
vari\^ancia da curva de loss em ${\sim}$50\% em rela\c{c}\~ao a
\emph{batch size}~32. Os modelos convergem entre 80--120 \'epocas;
configura\c{c}\~oes superiores convergem consistentemente em menos
\'epocas com menor loss de valida\c{c}\~ao, sustentando o
\emph{early stopping} como indicador confi\'avel de qualidade.
Todos os experimentos usam seed fixa (42) e opera\c{c}\~oes
determin\'{\i}sticas quando dispon\'{\i}veis.

\section{Distance-Aware Loss}

\subsection{O Problema}

A CE trata todos os erros igualmente. Fonologicamente, isso \'e
inadequado: substituir /\ipaE{}/ por /e/ (1 feature de diferen\c{c}a,
$d{=}0{,}10$) \'e qualitativamente diferente de substituir /a/ por /k/
(8+ features, $d{=}0{,}90$). Ambos recebem CE${}= 1{,}0$.

\subsection{A F\'ormula}

\begin{equation}
L = L_{CE} + \lambda \cdot d_{\text{PanPhon}}(\hat{y}_i, y_i)
    \cdot p_i^{(\hat{y}_i)}
\end{equation}

\noindent onde $d_{\text{PanPhon}}$ \'e a dist\^ancia calculada sobre 24
features articulat\'orias bin\'arias do PanPhon
\cite{mortensen2016panphon}, normalizada para $[0,1]$;
$p_i^{(\hat{y}_i)}$ \'e a probabilidade atribu\'{\i}da ao fonema predito
(fator de confian\c{c}a); e $\lambda$ \'e o peso do sinal fonol\'ogico.

A DA Loss adiciona penalidade proporcional a: (i) o quanto o modelo errou
(dist\^ancia) e (ii) o qu\~ao confiante estava no erro (probabilidade).

\textbf{Raz\~ao de design.} A CE monitora $p_\text{correto}$; a DA
monitora $p(\hat{y})$. Estas s\~ao independentes: quando o modelo erra
com confian\c{c}a, $p_\text{correto}$ \'e baixo enquanto $p(\hat{y})$
\'e alto. O fator de confian\c{c}a escala a penalidade articulat\'oria
com a certeza do modelo --- penalizando maximamente erros que s\~ao
simultaneamente confiantes e fonologicamente distantes.

\textbf{Sinal bounded.} A DA Loss \'e limitada superiormente por
$\lambda \times 1{,}0 \times 1{,}0 = 0{,}20$, enquanto a CE pode
atingir ${\sim}$16 no in\'{\i}cio do treinamento. A DA \'e efetiva
principalmente na zona de transi\c{c}\~ao (CE entre 0,3--1,5), quando
o modelo est\'a resolvendo ambiguidades entre fonemas candidatos.

\textbf{Exemplo num\'erico.} Considere dois erros com CE id\^entica
($p_\text{correto}{=}0{,}40$, CE${=}0{,}916$, $p(\hat{y}){=}0{,}45$),
sintetizados na Tabela~\ref{tab:example}.

\begin{table}[ht]
\centering
\caption{Exemplo num\'erico: dois erros com CE id\^entica mas
         dist\^ancia articulat\'oria distinta.}
\label{tab:example}
\begin{tabular}{llcccc}
\toprule
\textbf{Caso} & \textbf{Erro} & $d$ & \textbf{CE} & \textbf{DA} & $L$ total \\
\midrule
Near-miss      & e$\to$\ipaE{} & 0,10 & 0,916 & 0,009 & \textbf{0,925} \\
Catastr\'ofico & a$\to$k       & 0,90 & 0,916 & 0,081 & \textbf{0,997} \\
\bottomrule
\end{tabular}
\end{table}

A CE sozinha \'e id\^entica (0,916) em ambos os casos. A DA fornece
uma diferen\c{c}a de 9$\times$ na penalidade fonol\'ogica, permitindo
ao modelo distinguir um erro leve de um catastr\'ofico no mesmo passo
de gradiente.

\subsection{Sweep de $\lambda$}

A Tabela~\ref{tab:lambda} mostra o sweep de $\lambda$ com arquitetura
fixa (4,3M, sem separadores). O padr\~ao U-invertido motiva
$\lambda{=}0{,}20$: muito baixo deixa a CE indiferenciada; muito alto
compete com a CE.

\begin{table}[ht]
\centering
\caption{Sweep de $\lambda$, arquitetura 4,3M, sem separadores.}
\label{tab:lambda}
\begin{tabular}{cccc}
\toprule
$\lambda$ & PER & WER & Comportamento \\
\midrule
0,05 & 0,62\% & 5,36\% & Sinal fraco \\
0,10 & 0,63\% & 5,35\% & Melhoria moderada \\
\textbf{0,20} & \textbf{0,60\%} & \textbf{5,14\%} & \textbf{\'Otimo} \\
0,50 & 0,65\% & 5,57\% & Sobre-penaliza\c{c}\~ao \\
\bottomrule
\end{tabular}
\end{table}

\subsection{Dist\^ancias Customizadas}

O PanPhon atribui vetor zero a tokens n\~ao-fon\'eticos (separador
sil\'abico ``.'', marcador de acento ``\stressmark''), resultando em
$d(.{,}\text{\stressmark}) = 0{,}0$. \emph{Override}
p\'os-normaliza\c{c}\~ao seta $d{=}1{,}0$ para pares envolvendo
s\'{\i}mbolos estruturais, corrigindo o problema.

\section{Experimentos e Resultados}

\subsection{Progress\~ao dos Experimentos}

A Tabela~\ref{tab:experiments} sintetiza a progress\~ao experimental.

\begin{table}[ht]
\centering
\caption{Progress\~ao dos experimentos. Sep = separadores sil\'abicos.}
\label{tab:experiments}
\begin{tabular}{lcccccc}
\toprule
\textbf{Exp} & \textbf{Params} & \textbf{Loss} & \textbf{Sep}
  & \textbf{PER} & \textbf{WER} & \textbf{Insight} \\
\midrule
Exp0  & 4,3M  & CE       & n\~ao & 1,12\% & 9,37\% & Baseline \\
Exp1  & 4,3M  & CE       & n\~ao & 0,66\% & 5,65\% & $-$41\% PER \\
Exp5  & 9,7M  & CE       & n\~ao & 0,63\% & 5,38\% & Sweet spot \\
Exp7  & 4,3M  & DA 0,2   & n\~ao & 0,60\% & 5,14\% & $\lambda$ \'otimo \\
\textbf{Exp9}  & \textbf{9,7M}  & \textbf{DA 0,2}   & \textbf{n\~ao}
  & \textbf{0,58\%} & \textbf{4,96\%} & \textbf{Melhor WER} \\
Exp102 & 9,7M & CE       & sim   & 0,52\% & 5,79\% & Sep: PER$\downarrow$ WER$\uparrow$ \\
Exp103 & 9,7M & DA 0,2   & sim   & 0,53\% & 5,73\% & N\~ao-aditivos \\
\textbf{Exp104d} & \textbf{17,2M} & \textbf{DA+dist} & \textbf{sim}
  & \textbf{0,48\%} & \textbf{5,33\%} & \textbf{Ref.\ PER} \\
\bottomrule
\end{tabular}
\end{table}

\subsection{An\'alise Fatorial: Separadores $\times$ DA Loss}

Fatorial 2$\times$2 limpo (9,7M params, split 60/10/30), na
Tabela~\ref{tab:factorial}.

\begin{table}[ht]
\centering
\caption{Fatorial Separadores $\times$ DA Loss (PER / WER, 9,7M params).}
\label{tab:factorial}
\begin{tabular}{lcc}
\toprule
 & \textbf{CE} & \textbf{DA $\lambda{=}0{,}2$} \\
\midrule
Sem sep & 0,63\% / 5,38\% & \textbf{0,58\% / 4,96\%} \\
Com sep & 0,52\% / 5,79\% & 0,53\% / 5,73\% \\
\bottomrule
\end{tabular}
\end{table}

Separadores melhoram PER ($-$17\% a $-$20\%) mas pioram WER (+6\% a
+8\%). DA Loss melhora tanto PER quanto WER sem separadores. Com
separadores, o efeito da DA Loss \'e atenuado.

\subsection{Redistribui\c{c}\~ao Graduada dos Erros}

A Tabela~\ref{tab:classes} mostra a redistribui\c{c}\~ao das classes de
erro.

\begin{table}[ht]
\centering
\caption{Redistribui\c{c}\~ao das classes de erro com DA Loss.}
\label{tab:classes}
\begin{tabular}{llcccc}
\toprule
\textbf{Exp} & \textbf{T\'ecnica} & \textbf{PER}
  & \textbf{Cls~B} & \textbf{Cls~D} & \textbf{D/erros} \\
\midrule
Exp1 & CE baseline  & 0,66\% & 0,39\% & 0,54\% & 50,9\% \\
Exp6 & DA 0,1       & 0,63\% & 0,39\% & 0,47\% & 48,0\% \\
\textbf{Exp9} & \textbf{DA 0,2 (9,7M)} & \textbf{0,58\%}
  & \textbf{0,36\%} & \textbf{0,44\%} & \textbf{48,4\%} \\
\bottomrule
\end{tabular}
\end{table}

Classe~D (catastr\'ofica) cai de 0,54\% para 0,44\% ($-$19\%) enquanto
Classe~B se mant\'em. A DA Loss redistribui sistematicamente os erros.
Em aplica\c{c}\~oes TTS \emph{downstream}, erros Classe~B
(e.g., /e/$\leftrightarrow$/\ipaE{}/) s\~ao perceptualmente menos
salientes que erros Classe~D (e.g., vogal$\leftrightarrow$oclusiva);
valida\c{c}\~ao perceptual formal (MOS/ABX) permanece como trabalho
futuro.

\subsection{Evolu\c{c}\~ao das M\'etricas Graduadas}

Al\'em do PER cl\'assico, o \textbf{PER$_w$} (PER ponderado pela
dist\^ancia PanPhon) captura a \emph{qualidade} dos erros, n\~ao
apenas sua quantidade. A Tabela~\ref{tab:graduated} mostra a
evolu\c{c}\~ao.

\begin{table}[ht]
\centering
\caption{Evolu\c{c}\~ao das m\'etricas graduadas (sem separadores).}
\label{tab:graduated}
\begin{tabular}{llccc}
\toprule
\textbf{Exp} & \textbf{T\'ecnica} & \textbf{PER}
  & \textbf{PER$_w$} & \textbf{WER$_g$} \\
\midrule
Exp1 & CE baseline  & 0,66\% & 0,34\% & 0,69\% \\
Exp2 & CE, +capacidade & 0,60\% & 0,29\% & 0,62\% \\
Exp3 & CE + PanPhon$_T$ & 0,66\% & 0,33\% & 0,67\% \\
Exp6 & DA $\lambda{=}0{,}1$ & 0,63\% & 0,27\% & 0,60\% \\
\textbf{Exp9} & \textbf{DA $\lambda{=}0{,}2$ + 9,7M}
  & \textbf{0,58\%} & \textbf{0,27\%} & \textbf{0,58\%} \\
\bottomrule
\end{tabular}
\end{table}

A DA Loss reduz o PER$_w$ em $-$21\% (0,34\%$\to$0,27\%) em rela\c{c}\~ao
ao baseline, enquanto capacidade isolada (Exp2, 17,2M) reduz apenas
$-$15\%. Isto confirma que a DA Loss melhora a \emph{qualidade} dos erros
de forma desproporcional \`a redu\c{c}\~ao de quantidade.

\subsection{Compara\c{c}\~ao com LatPhon}

A refer\^encia mais pr\'oxima \'e o LatPhon \cite{chary2025latphon},
um Transformer multil\'{\i}ngue de 4 camadas (7,5M par\^ametros) que
reporta PER de 0,86\% (IC Wilson 95\% [0,56\%; 1,16\%]) em ${\sim}$500
palavras PT-BR do mesmo recurso lexical (ipa-dict). Ambos os sistemas
reportam ICs de Wilson; os intervalos n\~ao se sobrep\~oem --- o limite
superior do nosso sistema (0,51\%) fica abaixo do limite inferior do
LatPhon (0,56\%), diferen\c{c}a estatisticamente significativa a 95\%
de confian\c{c}a.

A diferen\c{c}a de 57$\times$ em escala de teste confere ICs
10$\times$ mais precisos \`a nossa avalia\c{c}\~ao ($\pm$0,03 p.p.\
vs.\ $\pm$0,30 p.p.). Diferen\c{c}as metodol\'ogicas impedem
compara\c{c}\~ao arquitetural direta: o LatPhon \'e um Transformer
multil\'{\i}ngue com NLL padr\~ao; nosso sistema \'e um BiLSTM
monol\'{\i}ngue com DA Loss. O resultado \'e consistente com a
hip\'otese de que desenho metodol\'ogico (fun\c{c}\~ao de custo +
protocolo de avalia\c{c}\~ao) pode compensar diferen\c{c}as
arquiteturais, sem estabelecer hierarquia universal entre fam\'{\i}lias
de modelo.

\section{An\'alise de Erros}

\subsection{Padr\~oes de Confus\~ao Dominantes}

A Tabela~\ref{tab:confusions} mostra as confus\~oes dominantes.
Mais de 60\% dos erros s\~ao neutraliza\c{c}\~oes voc\'alicas ---
substitui\c{c}\~oes entre vogais m\'edias abertas e fechadas
(/\ipaE{}/$\leftrightarrow$/e/,
/\ipaO{}/$\leftrightarrow$/o/). Estas refletem ambiguidade
fonol\'ogica genu\'{\i}na do PT-BR \cite{barbosa2004brazilian}: em
posi\c{c}\~ao \'atona, vogais m\'edias neutralizam-se.

\begin{table}[ht]
\centering
\caption{Principais confus\~oes fon\^emicas (Exp104d, 17,2M DA+dist).}
\label{tab:confusions}
\begin{tabular}{lrl}
\toprule
\textbf{Substitui\c{c}\~ao} & \textbf{Contagem} & \textbf{Mecanismo} \\
\midrule
\ipaE{}$\to$e & 255 & Neutraliza\c{c}\~ao \'atona \\
e$\to$\ipaE{} & 197 & Mesmo (dire\c{c}\~ao inversa) \\
\ipaO{}$\to$o & 131 & Neutraliza\c{c}\~ao posterior \\
i$\to$e      & 121 & Redu\c{c}\~ao voc\'alica \\
o$\to$\ipaO{} & 95 & Neutraliza\c{c}\~ao posterior \\
\bottomrule
\end{tabular}
\end{table}

A distribui\c{c}\~ao no corpus revela a causa estrutural: /e/ tem
raz\~ao 24,9:1 vs.\ /\ipaE{}/ em posi\c{c}\~ao pr\'e-t\^onica, mas
/\ipaE{}/ \'e dominante em posi\c{c}\~ao t\^onica (raz\~ao 0,33:1).
O modelo aprende o vi\'es pr\'e-t\^onico forte e o sobregeneraliza
para s\'{\i}labas t\^onicas --- artefato distribucional do corpus que
requer features pros\'odicos expl\'{\i}citos para ser resolvido.

\subsection{Regra Alof\^onica r-coda}

Durante a avalia\c{c}\~ao OOV, o modelo produziu /\ipaG{}/ em coda
antes de consoante vozeada (e.g., \emph{borboleta} $\to$
b~o~\ipaG{}~.~b~o\ldots). Auditoria do corpus revelou que o modelo
estava \textbf{correto}: a regra de assimila\c{c}\~ao de vozeamento do
PT-BR \cite{barbosa2004brazilian} exige /\ipaG{}/ antes de vozeada e
/x/ em coda final --- distribui\c{c}\~ao complementar perfeita no corpus
de 95.937 entradas.

\subsection{Avalia\c{c}\~ao de Generaliza\c{c}\~ao OOV}

Banco diagn\'ostico de 31 palavras em 6 categorias
(Tabela~\ref{tab:oov}).

\begin{table}[ht]
\centering
\caption{Avalia\c{c}\~ao de generaliza\c{c}\~ao OOV (Exp104d).}
\label{tab:oov}
\begin{tabular}{lccc}
\toprule
\textbf{Categoria} & \textbf{Corretas} & \textbf{Score} & \textbf{Insight} \\
\midrule
Generaliza\c{c}\~ao PT-BR & 6/9 (67\%) & 97\% & Near-misses \\
Consoantes Duplas   & 1/5 (20\%) & 81\% & Geminadas n\~ao vistas \\
Anglicismos         & 1/5 (20\%) & 71\% & Fonologia EN OOV \\
Chars OOV (k/w/y)   & 0/3 (0\%)  & 68\% & Falha esperada \\
\textbf{PT-BR Reais (OOV)} & \textbf{5/5 (100\%)} & \textbf{100\%}
  & \textbf{Generaliza\c{c}\~ao} \\
Controles           & 4/4 (100\%) & 100\% & Sanidade \\
\midrule
\textbf{Total} & \textbf{17/31 (55\%)} & & \\
\bottomrule
\end{tabular}
\end{table}

O resultado 5/5 em palavras reais in\'editas --- \emph{puxadinho},
\emph{abacatada}, \emph{zunido}, \emph{malcriado}, \emph{arrombado} ---
todas transcritas corretamente, \'e evid\^encia de que o modelo aprendeu
regras fonol\'ogicas produtivas do PT-BR (palataliza\c{c}\~ao,
redu\c{c}\~ao de coda, nasaliza\c{c}\~ao).

O banco de $N{=}31$ \'e uma \textbf{sonda diagn\'ostica}, n\~ao
substituto do test set principal de 28.782 palavras. Seu prop\'osito
\'e qualitativo: demonstrar que regras fonol\'ogicas espec\'{\i}ficas
transferem para palavras n\~ao vistas. Um conjunto adicional de 35
neolog\'{\i}smos (5 categorias: anglicismos, verbos emprestados,
neolog\'{\i}smos nativos, nomes pr\'oprios, palavras inventadas)
fornece evid\^encia complementar ao longo do mesmo eixo.

\textbf{Modos de falha definidos}: (1)~consoantes geminadas de
empr\'estimos italianos/ingleses, ausentes do corpus de treino;
(2)~anglicismos com fonologia inglesa, genuinamente OOV
fonol\'ogico; (3)~caracteres k, w, y --- OOV hard de charset.

\section{Discuss\~ao}

\textbf{Trade-off PER/WER com separadores.}
Separadores sil\'abicos melhoram PER ($-$17\% a $-$20\%) mas pioram WER
(+6\% a +8\%). O mecanismo \'e estrutural: cada token separador
mal-posicionado conta como erro de palavra inteira. A escolha depende da
aplica\c{c}\~ao: TTS prioriza PER; NLP/lookup prioriza WER.

\textbf{Limites da DA Loss.}
S\'{\i}mbolos n\~ao-fon\'eticos recebem vetor zero no PanPhon. O
\emph{override} corrige parcialmente, mas confus\~oes posicionais
persistem. Al\'em disso, DA \'e bounded por $\lambda \times 1{,}0 \times
1{,}0 = 0{,}20$, enquanto CE pode atingir ${\sim}$16 --- DA representa
$<$5\% do sinal quando o modelo est\'a muito errado, sendo efetiva na
zona de transi\c{c}\~ao.

\textbf{Memoriza\c{c}\~ao vs.\ aprendizado.}
O split 60/10/30 prioriza test set grande sobre treino m\'aximo.
Com 17,2M par\^ametros e 57.562 palavras de treino, o modelo
\emph{poderia} memorizar --- uma tabela hash seria mais compacta.
Contudo, o acerto de 100\% em palavras reais in\'editas (com
palataliza\c{c}\~ao d+i$\to$d\textipa{Z}, redu\c{c}\~ao de coda
l$\to$w, nasaliza\c{c}\~ao om$\to$\~{o} corretos) \'e consistente
com aquisi\c{c}\~ao de regras fonol\'ogicas produtivas.
Al\'em disso, reduzir os dados de treino de 60\% para 50\% (Exp105)
degrada o PER em apenas 10\% relativo (0,49\%$\to$0,54\%), sugerindo
que o corpus est\'a bem acima do limiar de satura\c{c}\~ao para esta
arquitetura.

\textbf{Intera\c{c}\~ao de capacidade.}
Em 17,2M par\^ametros \emph{sem} separadores (Exp10), a DA Loss n\~ao
melhora o PER em rela\c{c}\~ao \`a CE (0,61\% vs.\ 0,60\%).
Contudo, em 17,2M \emph{com} separadores e dist\^ancias corrigidas
(Exp104d), a DA atinge o melhor PER observado (0,48\%). Isto refuta a
hip\'otese simples de que ``modelos grandes memorizam e a DA
atrapalha'' --- a vari\'avel relevante \'e a \emph{intera\c{c}\~ao}
entre capacidade, representa\c{c}\~ao estrutural e corre\c{c}\~ao de
dist\^ancia.

\subsection{Trabalhos Futuros}

Tr\^es dire\c{c}\~oes promissoras:
(1)~\textbf{Restri\c{c}\~oes fonot\'aticas}: O PT-BR restringe codas
a 4 consoantes (/s/, /\textipa{R}/, /l/, /N/) e ${\sim}$78\% das
s\'{\i}labas s\~ao abertas (CV). Um aut\^omato de estados finitos
codificando a estrutura sil\'abica (ONSET$\to$NUCLEUS$\to$CODA)
poderia reduzir erros de posicionamento de separadores.
(2)~\textbf{Espa\c{c}o articulat\'orio cont\'{\i}nuo 7D}: Substituir
os 24 features bin\'arios do PanPhon por um espa\c{c}o vocal cont\'{\i}nuo
de 7 dimens\~oes atribuiria coordenadas distintas aos tokens
estruturais, eliminando a necessidade de \emph{override} de
dist\^ancias.
(3)~\textbf{Contexto morfossint\'atico}: Hom\'ografos heter\'ofonos
(e.g., \emph{jogo} substantivo /\ipaO{}/ vs.\ verbo /o/)
requerem contexto sentencial al\'em do G2P no n\'{\i}vel de
palavra.

\section{Limita\c{c}\~oes}

\begin{enumerate}
\item \textbf{Conjunto OOV pequeno}: O banco de 31 palavras \'e sonda
  diagn\'ostica, n\~ao avalia\c{c}\~ao exaustiva.
\item \textbf{Monolinguismo}: Sistema treinado e avaliado apenas em PT-BR.
\item \textbf{Valida\c{c}\~ao perceptual ausente}: Hip\'otese de que erros
  pr\'oximos s\~ao menos sal\'{\i}veis em TTS n\~ao validada (MOS/ABX).
\item \textbf{Corpus \'unico}: Generaliza\c{c}\~ao para outros corpora
  PT-BR n\~ao testada.
\item \textbf{Arquitetura \'unica}: DA Loss avaliada apenas com BiLSTM.
\end{enumerate}

\section{Declara\c{c}\~ao de \'Etica}

Este trabalho utiliza um dicion\'ario fon\'etico publicamente
dispon\'{\i}vel sem informa\c{c}\~oes pessoais identific\'aveis.
N\~ao foram realizados experimentos com participantes humanos.
Os modelos treinados convertem texto em representa\c{c}\~ao fon\'etica e
n\~ao geram conte\'udo potencialmente prejudicial.

\section{Uso de IA Generativa}

Ferramentas de IA generativa (Claude, Anthropic) foram utilizadas
exclusivamente como assist\^encia editorial: revis\~ao gramatical e
ortogr\'afica, pesquisa de refer\^encias bibliogr\'aficas, gest\~ao de
cita\c{c}\~oes, verifica\c{c}\~ao de consist\^encia terminol\'ogica entre
se\c{c}\~oes e formata\c{c}\~ao \LaTeX. A ferramenta n\~ao gerou texto
novo, hip\'oteses, figuras, c\'odigo nem interpreta\c{c}\~oes de dados.
Todo o conte\'udo cient\'{\i}fico --- design experimental,
implementa\c{c}\~ao, execu\c{c}\~ao dos experimentos, an\'alise dos
resultados e conclus\~oes --- \'e de autoria exclusiva dos autores.

\bibliographystyle{sbc}
\bibliography{stil_refs}

\end{document}
